\documentclass{article}

\usepackage[english]{babel}
\usepackage[letterpaper,top=2cm,bottom=2cm,left=3cm,right=3cm,marginparwidth=1.75cm]{geometry}
\usepackage{amsmath}
\usepackage{amssymb}
\usepackage{graphicx}
\usepackage[colorlinks=true, allcolors=blue]{hyperref}
\usepackage[round]{natbib}

\DeclareMathOperator{\sgn}{sgn}
\newcommand{\RE}{\mathrm{Re}}
\newcommand{\IM}{\mathrm{Im}}
\newcommand{\HT}{\mathcal{H}}
\newcommand{\Hfp}{\operatorname{Hfp}}
\counterwithin*{equation}{section}
\renewcommand\theequation{\thesection.\arabic{equation}}

\title{The Fractional Laplacian: Rigorous Construction and Physical Origin\\
  of the Half-Order Diffusion Operator\\
  \large Supplementary Annex}
\author{Navarrete, Iv\'an \and Olive, Jean-Arthur}

\begin{document}
\maketitle

\begin{abstract}
Section 2 of the main text reduces the governing equation for inter-rifting
deformation to $\partial_t\hat u = -a|k|\hat u$ in Fourier space, and
identifies $|k|$ as the symbol of the half-Laplacian $(-\Delta)^{1/2}$. This
annex grounds that identification rigorously, following the definitions and
equivalences surveyed by \citet{daoud2022fractional}. The central result is
that the operator order $s=1/2$ is not a modeling choice: the
Caffarelli--Silvestre extension \citep{caffarelli2007extension} that
represents $(-\Delta)^s$ as a Dirichlet-to-Neumann map degenerates, exactly
at $s=1/2$, to the classical unweighted harmonic extension problem, the same
half-space potential problem underlying the Newtonian viscous half-space
coupling of \citet{savage1978asthenosphere}. We give short proofs of the
results that are load-bearing for the model's own construction, cite the
rest, and synthesize the mathematics for a structural-geology reading of
the model.
\end{abstract}

\tableofcontents

\section{The kernel is the fractional Laplacian}
\label{sec:kernel}

\subsection{Recap of the operator}

Section 2.1--2.2 of the main text derives the governing equation for the
surface displacement $u(x,t)$ of an elastic plate over a Newtonian viscous
half-space, and reduces it in Fourier space to
\begin{equation}
\frac{\partial\hat u}{\partial t} = -a|k|\hat u,
\label{frac:governing}
\end{equation}
with $a = EH/[2\eta(1-\nu^2)]$. The wavenumber symbol $|k|$ is identified
there, by reference to the classical fractional-diffusion literature
\citep{metzler2000random,samko1993fractional}, as the Fourier symbol of the
half-Laplacian $(-\Delta)^{1/2}$. This annex does not re-derive the physical
setup; it establishes, from the definitions surveyed by
\citet{daoud2022fractional}, which operator this is and why its
order is forced to be $s=\tfrac12$ by the coupling's own geometry,
not a relabeling of an already-known kernel.

\subsection{The normalizing constant at $N=1$, $s=\tfrac12$}

\citet{daoud2022fractional} give the fractional Laplacian on $\mathbb{R}^N$
as the hypersingular principal-value integral. Here and for the rest of
this annex, $v$ denotes a generic test function on which the operator
acts, following \citeauthor{daoud2022fractional}'s own notation; it is
unrelated to the surface velocity $\partial u/\partial t$ or to the main
text's $v^{\mathrm{STEP}}$, which never appears in this abstract role.
\begin{equation}
(-\Delta)^s v(x) = C(N,s)\;\mathrm{P.V.}\!\int_{\mathbb{R}^N}
  \frac{v(x)-v(y)}{\|x-y\|^{N+2s}}\,dy,
\qquad
C(N,s) = \frac{s\,4^s}{\pi^{N/2}}\,\frac{\Gamma(s+\tfrac{N}{2})}{\Gamma(1-s)}.
\label{frac:singint}
\end{equation}
At $N=1$, $s=\tfrac12$,
\begin{equation}
C\!\left(1,\tfrac12\right)
= \frac{\tfrac12\cdot 4^{1/2}}{\pi^{1/2}}\cdot\frac{\Gamma(1)}{\Gamma(\tfrac12)}
= \frac{1}{\sqrt\pi}\cdot\frac{1}{\sqrt\pi}
= \boxed{\frac{1}{\pi}},
\label{frac:C1half}
\end{equation}
using $\Gamma(1)=1$ and $\Gamma(\tfrac12)=\sqrt\pi$. Confirmed numerically in
\texttt{verify\_fractional\_kernel.jl}. So, for $N=1$ and $s=\tfrac12$,
(\ref{frac:singint}) reads
\begin{equation}
(-\Delta)^{1/2}v(x) = \frac{1}{\pi}\;\mathrm{P.V.}\!\int_{-\infty}^{\infty}
  \frac{v(x)-v(y)}{(x-y)^2}\,dy.
\label{frac:half}
\end{equation}
The factor of $\pi$ in the denominator of Savage's kernel
(main text eq.~3, $2\eta/[\pi(x-x')^2]$) is not a rheological constant: it
is exactly the normalizing constant that makes the singular integral
definition of the half-Laplacian equal its Fourier-multiplier definition.

\subsection{Reconciling the two principal values}
\label{sec:hadamard}

There is a subtlety that must be handled carefully rather than glossed
over. The Stokes kernel convolution (main text eq.~3) is
\begin{equation}
\sigma_{xz}^{\mathrm{BOT}}(x,t) = \int_{-\infty}^{\infty}
  \frac{2\eta}{\pi(x-x')^2}\,\frac{\partial u}{\partial t}(x',t)\,dx',
\label{frac:stokes}
\end{equation}
described, following \citet{savage1978asthenosphere}, as a Cauchy principal
value acting directly on $\partial u/\partial t$. Definition~(\ref{frac:half})
instead takes the principal value of the \emph{difference} $v(x)-v(y)$, not
of $v(y)$ alone. These are not the same regularization, and the difference is
not cosmetic: a genuine symmetric-truncation principal value of
$\int v(y)/(x-y)^2\,dy$ diverges for any $v(x)\neq0$, because $1/(x-y)^2$ is
positive on both sides of the singularity and the truncated integral grows
as $2v(x)/\varepsilon$ with no cancellation. What~(\ref{frac:stokes}) calls a
principal value is, precisely, the Hadamard finite part,
\begin{equation}
\Hfp\!\int_{-\infty}^{\infty}\frac{v(y)}{(x-y)^2}\,dy
:= \lim_{\varepsilon\to0^+}\left[
    \int_{|y-x|\ge\varepsilon}\frac{v(y)}{(x-y)^2}\,dy - \frac{2v(x)}{\varepsilon}
  \right],
\label{frac:hfp}
\end{equation}
a regularization standard in dislocation and crack theory for this
order-2 kernel, loosely called a "principal value" in that same literature.

The two regularizations are related by a clean sign identity. Expanding the
principal value of the difference form,
\begin{align}
\mathrm{P.V.}\!\int\frac{v(x)-v(y)}{(x-y)^2}\,dy
&= \lim_{\varepsilon\to0^+}\left[
     v(x)\!\int_{|y-x|\ge\varepsilon}\!\frac{dy}{(x-y)^2}
     - \int_{|y-x|\ge\varepsilon}\!\frac{v(y)}{(x-y)^2}\,dy
   \right] \nonumber\\
&= \lim_{\varepsilon\to0^+}\left[
     \frac{2v(x)}{\varepsilon}
     - \int_{|y-x|\ge\varepsilon}\!\frac{v(y)}{(x-y)^2}\,dy
   \right]
 = -\Hfp\!\int\frac{v(y)}{(x-y)^2}\,dy,
\label{frac:sign}
\end{align}
using $\int_{|z|\ge\varepsilon}dz/z^2 = 2/\varepsilon$ exactly, for every
$\varepsilon>0$. Confirmed numerically in
\texttt{verify\_fractional\_kernel.jl} for a Gaussian test function.
Combining~(\ref{frac:half}) and~(\ref{frac:sign}),
\begin{equation}
\Hfp\!\int\frac{v(y)}{(x-y)^2}\,dy = -\pi\,(-\Delta)^{1/2}v(x).
\label{frac:hfp_final}
\end{equation}

Applying~(\ref{frac:hfp_final}) to $v=\partial u/\partial t$ in the Stokes
kernel~(\ref{frac:stokes}),
\begin{equation}
\sigma_{xz}^{\mathrm{BOT}}(x,t)
= \frac{2\eta}{\pi}\,\Hfp\!\int\frac{\partial u/\partial t\,(x')}{(x-x')^2}\,dx'
= \frac{2\eta}{\pi}\left[-\pi\,(-\Delta)^{1/2}\frac{\partial u}{\partial t}\right]
= -2\eta\,(-\Delta)^{1/2}\frac{\partial u}{\partial t}.
\label{frac:kernel-identity}
\end{equation}
In Fourier space, $(-\Delta)^{1/2}$ has symbol $|k|$, so
$\mathcal{F}\{\sigma_{xz}^{\mathrm{BOT}}\} = -2\eta|k|\,\partial\hat u/\partial t$,
the classical result quoted in the main text (eq.~5) and in
\citet{savage1978asthenosphere}. The Stokes kernel is therefore $-2\eta$
times the half-Laplacian of the surface velocity field, term for term and
with the correct sign. This is not a loose spectral analogy: identifying
the regularization used in the geophysics literature with the Hadamard
finite part implicit in the singular-integral definition of
$(-\Delta)^{1/2}$ shows the two are the same operator.

This identity holds for the Stokes term alone, and it is naturally a
statement about velocity rather than displacement: a viscous material
resists the \emph{rate} of deformation, not the deformation itself, so the
basal traction depends on $\partial u/\partial t$, unlike the elastic term
(main text eq.~1, Hooke's law), which depends on $u$ directly. Combining
the two in the full governing equation (main text eq.~4) and
Fourier-transforming is what closes the loop back to~(\ref{frac:governing}).
Because $(-\Delta)^{1/2}$ is a purely spatial, time-independent operator,
it commutes with $\partial/\partial t$:
$(-\Delta)^{1/2}(\partial u/\partial t) = \partial_t\big[(-\Delta)^{1/2}u\big]$.
The combined equation therefore rearranges into an equation purely in $u$,
$\partial_t\hat u = -a|k|\hat u$: the same half-Laplacian that maps
velocity to traction in the viscous term also governs the time evolution
of displacement itself, once the elastic restoring term is folded in. From
here on, this annex works with that displacement-evolution form directly
(the initial-value problem for $u(x,t)$ given $u(x,0)$), since that is the
equation whose solution is the object of physical interest; the velocity
field never needs to reappear as a separate unknown.

\section{The correct function space: $H^{1/2}(\mathbb{R})$}
\label{sec:sobolev}

\citet{daoud2022fractional} identify $H^s(\mathbb{R}^N)$, equipped with the
Gagliardo seminorm
\begin{equation}
[v]_{H^s(\mathbb{R}^N)}^2 := \int_{\mathbb{R}^N}\!\int_{\mathbb{R}^N}
  \frac{|v(x)-v(y)|^2}{\|x-y\|^{N+2s}}\,dx\,dy,
\label{frac:gagliardo}
\end{equation}
as the natural energy space for $(-\Delta)^s$: their Proposition~2.6 relates
the seminorm directly to the operator, $[v]_{H^s}^2 = 2\,C(N,s)^{-1}
\|(-\Delta)^{s/2}v\|_{L^2}^2$ (see also \citealp{dinezza2012hitchhikers} for
the full development of these spaces). For the model's own construction,
$N=1$, $s=\tfrac12$, so the relevant space is $H^{1/2}(\mathbb{R})$,
\begin{equation}
[v]_{H^{1/2}(\mathbb{R})}^2 = \int_{-\infty}^{\infty}\!\int_{-\infty}^{\infty}
  \frac{|v(x)-v(y)|^2}{(x-y)^2}\,dx\,dy.
\label{frac:gagliardo12}
\end{equation}

\subsection{The dike-opening initial condition is not in $H^{1/2}(\mathbb{R})$}
\label{sec:heaviside}

The initial condition for a single dike intrusion is a pure step. Following
the main text's own convention (eq.~11), $u(x,0^+) = b\,H(-x)$, where
$H(-x)$ is the Heaviside function ($1$ for $x<0$, $0$ for $x>0$); we write
it as plain $H$ throughout this annex, reserving the calligraphic $\HT$ for
the Hilbert transform introduced in Section~\ref{sec:extension} so the two
are never confused.

The Gagliardo seminorm is quadratic in its argument: for any constant $c$
and function $w$, $(cw)(x)-(cw)(y) = c\,(w(x)-w(y))$, so $c^2$ factors out
of the double integral in (\ref{frac:gagliardo12}) unchanged and
$[cw]_{H^{1/2}(\mathbb{R})}^2 = c^2\,[w]_{H^{1/2}(\mathbb{R})}^2$. It is
therefore enough to compute the seminorm of the unit-amplitude step
$v:=H(-\cdot)$ (the function $x\mapsto H(-x)$, with $b=1$) and rescale by
$b^2$ at the end, rather than carrying the physical amplitude $b$ through
every line of the calculation. Substituting $v=H(-\cdot)$ into (\ref{frac:gagliardo12}):
for $x>0$, $y<0$ (or vice versa), $v(x)=0$ and $v(y)=1$, so
$|v(x)-v(y)|^2=1$; for $x,y$ of the same sign, the integrand vanishes. So
\begin{equation}
[H(-\cdot)]_{H^{1/2}(\mathbb{R})}^2
= 2\int_0^\infty\!\int_0^\infty \frac{dt\,dx}{(x+t)^2},
\qquad y=-t,\ t>0,
\label{frac:sgn_seminorm}
\end{equation}
using the symmetry of the integrand under $(x,y)\mapsto(y,x)$ to double the
$x>0,y<0$ contribution. The inner integral is elementary and exact for every
$x>0$:
\begin{equation}
\int_0^\infty \frac{dt}{(x+t)^2} = \frac{1}{x}.
\end{equation}
The outer integral $\int_0^\infty dx/x$ diverges logarithmically as
$x\to0^+$ (confirmed numerically in
\texttt{verify\_sobolev\_divergence.jl}: the truncated integral
$I(\delta)=\int_\delta^\infty dx/x$ tracks $-\log\delta$ to within a fixed
offset across four orders of magnitude in $\delta$). Since $u(x,0^+) =
b\,H(-x)$ differs from $H(-\cdot)$ only by the constant scale factor $b$,
\begin{equation}
[u(\cdot,0^+)]_{H^{1/2}(\mathbb{R})}^2 = b^2\,[H(-\cdot)]_{H^{1/2}(\mathbb{R})}^2 = +\infty.
\label{frac:sobolev-div}
\end{equation}
The pure step is the textbook borderline case: $H^{1/2}(\mathbb{R})$ is the
critical exponent in one dimension, the threshold at which the fractional
Sobolev embedding fails to control jump discontinuities by only a
logarithm \citep{dinezza2012hitchhikers}. This is not a defect of the model.
It is a precise statement of a fact the model already handles correctly by
construction: the solution is not required to remain in the natural energy
space at $t=0$, because the evolution operator smooths it instantaneously
for any $t>0$, which we establish rigorously in Section~\ref{sec:semigroup}.

This raises a natural question: if $u(\cdot,0^+)\notin H^{1/2}(\mathbb{R})$,
where does it actually live, and does the failure matter? The two parts of
the question have different answers depending on whether "living" is asked
locally or globally. Locally, near the jump alone, the fractional Sobolev
embedding fails only at the critical exponent: repeating the
computation above with a generic exponent $s$ in place of $\tfrac12$, the
inner integral becomes $\int_0^\infty dt/(x+t)^{1+2s} = x^{-2s}/(2s)$, whose
outer integral over $x$ near $0$ converges for every $s<\tfrac12$ and
diverges (worse than logarithmically) for every $s>\tfrac12$; $s=\tfrac12$
is the one exponent at which it diverges only logarithmically. So
$u(\cdot,0^+)\in H^s_{\mathrm{loc}}(\mathbb{R})$ for every $s<\tfrac12$, and
fails only at, and above, the critical exponent. Globally, the more mundane
fact that $H(-x)$ never decays as $x\to\pm\infty$ keeps it out of
$L^2(\mathbb{R})$, and hence out of $H^s(\mathbb{R})$ for any $s$, for a
reason that has nothing to do with the jump; the natural global home for a
bounded function with a single finite jump and no other structure is the
space of bounded variation, $BV(\mathbb{R})$. So,
\begin{equation*}
u(\cdot,0^+)\in BV(\mathbb{R})\cap\!\!\bigcap_{s<1/2}\!\!H^s_{\mathrm{loc}}(\mathbb{R}),
\qquad u(\cdot,0^+)\notin H^{1/2}(\mathbb{R}),
\end{equation*}
the classical statement that $BV(\mathbb{R})$ fails to embed into the
critical Sobolev space $H^{1/2}(\mathbb{R})$ by exactly a logarithm
\citep{dinezza2012hitchhikers}.

Whether this matters depends on which route is used to solve the equation.
It does not matter for the route this annex takes: the semigroup formula
of Section~\ref{sec:semigroup} is defined for any tempered distribution as
initial data, $BV(\mathbb{R})$ included, with no need for the initial
condition to sit in the equation's own energy space, and the semigroup's
smoothing repairs whatever regularity the initial data lacks the instant
$t$ becomes positive. It would matter for a different route: classical
energy or variational methods for parabolic-type equations (Galerkin
approximation, weak formulations built directly on the natural energy
space) generally require the initial data to already sit in that energy
space to set up the weak problem cleanly, and a rigorous treatment via that
route would need to work around this initial condition explicitly, for
instance by posing the problem in $BV(\mathbb{R})$ directly or by falling
back on the semigroup formula used here.

\section{Construction via the semigroup}
\label{sec:semigroup}

This section proceeds in four steps. Section~\ref{sec:semigroup-construction}
constructs the model's solution directly and identifies it as a genuine
semigroup. Section~\ref{sec:semigroup-subordination} gives an equivalent,
second construction that puts the heat-semigroup function $w$
of~(\ref{frac:heatsemi}) to explicit use, rather than leaving it as an
unused ingredient of a definition. Section~\ref{sec:semigroup-analyticity}
shows what the semigroup does at high spatial frequency: it smooths, with
an explicit, growing radius of analyticity. Section~\ref{sec:semigroup-k0}
shows what it does not touch, at zero frequency, and why that is harmless.

\subsection{Construction of the solution}
\label{sec:semigroup-construction}

Equation~(\ref{frac:governing}) is not a postulate about an abstract
operator: it is the model's own momentum balance, already established in
Section~\ref{sec:kernel}, Hooke's law for the elastic plate combined with
the Stokes-kernel coupling to the viscous half-space, Fourier-transformed.
Solving it needs nothing beyond ordinary linear ODE theory. For each fixed
wavenumber $k$, $\partial_t\hat u = -a|k|\hat u$ integrates directly to
\begin{equation}
\hat u(k,t) = \hat u(k,0)\,e^{-a|k|t}.
\label{frac:semigroup}
\end{equation}
This is the model's solution, in full; nothing that follows in this
section adds to it. The family of multipliers $S(t):k\mapsto e^{-a|k|t}$,
$t\ge0$, satisfies the semigroup law exactly,
\begin{equation}
S(t_1)S(t_2) = e^{-a|k|t_1}e^{-a|k|t_2} = e^{-a|k|(t_1+t_2)} = S(t_1+t_2),
\qquad S(0)=1,
\label{frac:semigrouplaw}
\end{equation}
identifying it as the semigroup generated by $-a(-\Delta)^{1/2}$: solving
the physical equation and identifying the operator it contains are two
separate facts, and~(\ref{frac:semigrouplaw}) records the second, not a
further step in the first.

That the operator appearing in~(\ref{frac:governing}) is the same
$(-\Delta)^{1/2}$ characterized throughout this annex, rather than merely
something that behaves like it in Fourier space, is what
Section~\ref{sec:kernel} already proves directly for this specific kernel.
A further, independent confirmation comes from the general theory of the
operator itself, unrelated to the model or to $u$.
\citet{daoud2022fractional} give a definition of $(-\Delta)^s$, for any
generic function $v$, through the heat semigroup: for $w(x,\tau)$ solving
the classical heat equation $\partial_\tau w = \partial_x^2 w$ with
$w(x,0)=v(x)$,
\begin{equation}
(-\Delta)^s v(x) = \frac{s}{\Gamma(1-s)}\int_0^\infty
  \big(v(x)-w(x,\tau)\big)\,\frac{d\tau}{\tau^{1+s}}.
\label{frac:heatsemi}
\end{equation}
Here $\tau$ is an internal, dummy diffusion time used only to define the
operator at a single instant; it has no relation to the model's physical
evolution time $t$, and $v$ is any test function, not the model's own $u$.
Equation~(\ref{frac:heatsemi}) is pure operator theory: it holds regardless
of whether $(-\Delta)^{1/2}$ ever appears in a physical equation, and it
plays no role in solving~(\ref{frac:governing}), which was already solved
above with no reference to $w$ at all. Its purpose here is narrower and
different: to complete the catalogue of equivalent constructions of
$(-\Delta)^{1/2}$ (alongside the singular integral of
Section~\ref{sec:kernel} and the Fourier multiplier), and to set up a
second, genuinely useful role for $w$ in the model's own solution, taken up
next.

\subsection{An equivalent view: Bochner subordination and the role of $w$}
\label{sec:semigroup-subordination}

There is a second, equally valid route
to~(\ref{frac:semigroup}) that puts $w$ to direct use, rather than leaving
it stranded in~(\ref{frac:heatsemi}). Bochner's subordination formula
writes the semigroup generated by $-a(-\Delta)^{1/2}$ as a weighted average,
over all internal diffusion times $\tau$, of the ordinary heat semigroup:
\begin{equation}
u(x,t) = \int_0^\infty w(x,\tau)\,f_{at}(\tau)\,d\tau,
\qquad
f_{at}(\tau) = \frac{at}{2\sqrt\pi}\,\tau^{-3/2}\,e^{-(at)^2/(4\tau)},
\label{frac:subordination}
\end{equation}
where $w(x,\tau)$ is the ordinary (Fickian) heat-equation solution with the
model's own initial condition, $w(x,0)=u(x,0^+)=b\,H(-x)$, and $f_{at}$ is
the classical density of the $\tfrac12$-stable subordinator at
"time" $at$: the same $w$ as in~(\ref{frac:heatsemi}), now integrated
against an explicit weight instead of differenced against a single value
of $v$, so that~(\ref{frac:heatsemi}) and~(\ref{frac:subordination}) are
the differential and integrated statements of the same subordination
principle.

For the model's own initial condition, $w$ has the classical closed form
\begin{equation}
w(x,\tau) = \frac{b}{2}\\mathrm{erfc}\!\left(\frac{x}{2\sqrt\tau}\right),
\label{frac:wclosed}
\end{equation}
obtained by convolving $u(x,0^+)=b\,H(-x)$ with the ordinary heat kernel;
this is the Foulger diffusion profile the main text rejects as the
physical model (§4.3, erfc tails rather than Lorentzian). It reappears
here purely as a mathematical ingredient: the correct, non-Fickian solution
$u(x,t)$ is not $w(x,t)$ itself, but the average of $w(x,\cdot)$ over all
diffusion times $\tau$, weighted by the subordinator density $f_{at}$.

Two independent checks confirm~(\ref{frac:subordination}), at two different
levels. The first is general, at the level of the operator itself, and
does not reference the model's specific initial condition. The ordinary
heat semigroup has Fourier multiplier $e^{-k^2\tau}$ (a direct consequence
of $\partial_\tau w=\partial_x^2w$); averaging it against $f_{at}(\tau)$ is
exactly a Laplace transform of $f_{at}$ in $\tau$, evaluated at $p=k^2$:
\begin{equation}
\int_0^\infty e^{-k^2\tau}f_{at}(\tau)\,d\tau = e^{-at|k|},
\label{frac:laplacecheck}
\end{equation}
confirmed to machine precision for $k=0.3,1,2.5,5$. The right-hand side is
exactly the multiplier of~(\ref{frac:semigroup}), so this identity alone
already guarantees that~(\ref{frac:subordination})
and~(\ref{frac:semigroup}) agree for \emph{any} initial data $v$, not only
$u(x,0^+)$.

The second check applies that general fact to the model's own concrete
quantities. It is not logically required once~(\ref{frac:laplacecheck})
is established, but it ties the abstract operator identity back to the
specific numbers the model uses. Substituting the closed
form~(\ref{frac:wclosed}) into~(\ref{frac:subordination}) gives the
explicit, one-dimensional integral
\begin{equation}
u(x,t) = \int_0^\infty \frac{b}{2}\,
  \mathrm{erfc}\!\left(\frac{x}{2\sqrt\tau}\right)
  \frac{at}{2\sqrt\pi}\,\tau^{-3/2}\,e^{-(at)^2/(4\tau)}\,d\tau,
\label{frac:erfccheck}
\end{equation}
with no unresolved symbols left: every quantity on the right is a known
function of $\tau$ for fixed $x$, $a$, $b$, $t$. Evaluating~(\ref{frac:erfccheck})
by numerical quadrature reproduces the main text's own closed form
$u(x,t)=\tfrac{b}{2}-\tfrac{b}{\pi}\arctan(x/(at))$ to $10^{-9}$ across
$x\in[-3,3]$ (quadrature tolerance).

The subordination view is not needed for anything that follows, but it is
what makes the smoothing established next intuitive rather than merely
algebraic, since each $w(x,\tau)$, $\tau>0$, is already smooth by ordinary
heat-kernel diffusion, and an average of smooth functions is smooth.

\subsection{High-frequency damping and analyticity}
\label{sec:semigroup-analyticity}

Return to the direct construction~(\ref{frac:semigroup}). For fixed $t>0$
and any integer $n\ge0$, maximizing $k^ne^{-akt}$ over $k>0$ gives the
sharp bound
\begin{equation}
\sup_{k\in\mathbb{R}} |k|^n e^{-a|k|t} = \left(\frac{n}{ate}\right)^n < \infty,
\label{frac:supbound}
\end{equation}
confirmed numerically to $10^{-10}$ precision for $n=1,2,3,5$. Because this
supremum is finite for every $n$, $|k|^ne^{-a|k|t}\to0$ as $|k|\to\infty$
for every $n$: the multiplier decays faster than any power of $|k|$,
uniformly in $k$.

This decay controls the local smoothness of $u(\cdot,t)$ directly, and
makes it precise enough to quote a radius, not just an adjective. Extending
$x\mapsto x+i\sigma$ in~(\ref{frac:semigroup}) and inverting,
\begin{equation}
u(x+i\sigma,t) = \frac{1}{2\pi}\int_{-\infty}^{\infty}
  \hat u(k,0)\,e^{-a|k|t-k\sigma}\,e^{ikx}\,dk,
\label{frac:strip}
\end{equation}
the exponent $-a|k|t-k\sigma$ tends to $-\infty$ as $|k|\to\infty$ exactly
when $|\sigma|<at$: it grows as $-k(at+\sigma)$ for $k\to+\infty$ and as
$-|k|(at-\sigma)$ for $k\to-\infty$, both negative if and only if
$-at<\sigma<at$. So~(\ref{frac:strip}) converges, and $u(\cdot,t)$ extends
holomorphically, throughout the strip
\begin{equation}
|\mathrm{Im}\,z| < at,
\label{frac:stripwidth}
\end{equation}
confirmed numerically: truncating~(\ref{frac:strip}) at $|k|\le K$, the
truncated integral converges to a finite limit as $K\to\infty$ for every
$|\sigma|<at$ (e.g.\ stabilizing near $10.5$ for $\sigma/at=0.9$), grows
linearly in $K$ exactly at $\sigma=at$, and diverges super-exponentially
for $|\sigma|>at$ (already $10^{43}$ at $K=200$ for $\sigma/at=1.5$).
A function holomorphic throughout a strip around the real axis is, in
particular, real analytic on the real axis; (\ref{frac:stripwidth}) makes
this quantitative: the radius of convergence of $u(\cdot,t)$'s Taylor
series at any point $x\in\mathbb{R}$ is at least $at$, growing linearly in
time, no matter how singular $u(\cdot,0)$ is. The dike intrusion initial
condition $u(x,0^+)=b\,H(-x)$ has a jump at $x=0$ and $t=0$, and
that discontinuity is erased completely (radius of convergence $at>0$
everywhere) the instant $t$ becomes positive. This resolves
the tension flagged in Section~\ref{sec:heaviside}: $u(x,0^+)=b\,H(-x)$
fails to belong to the natural energy space $H^{1/2}(\mathbb{R})$, yet
$u(\cdot,t)$ is real analytic in $x$ for every $t>0$. The same conclusion
applies to the closed-form step response given in the main text (§2.3): its
printed formula is a single-branch simplification valid for $x>0$, but the
solution it represents, correctly extended to all $x$, is the real analytic
function~(\ref{frac:stripwidth}) predicts.

\subsection{What survives at $k=0$}
\label{sec:semigroup-k0}

Everything in Section~\ref{sec:semigroup-analyticity} concerned the limit
$|k|\to\infty$; the opposite limit, $k\to0$, is where the multiplier
$e^{-a|k|t}$ is least effective; it equals $1$ there for every $t$, so
whatever singularity $\hat u(k,0)$ carries at $k=0$ is carried forward
unchanged, for all time. This is not a gap in the argument above: it is a
separate question, about a different part of the spectrum, with its own
separate (and equally explicit) answer. The main text's own Fourier
transform of the initial condition (eq.~12) is
$\hat u(k,0)=b\pi\delta(k)+ib/k$, and $e^{-a|k|t}=1$ at $k=0$ for every $t$,
so neither term is ever damped there. The Dirac term $b\pi\delta(k)$ alone
is enough to show
$\hat u(\cdot,t)$ is never Schwartz class, for any $t\ge0$: a genuine delta
distribution is not even a locally bounded function, let alone
rapidly decaying. It corresponds, on the spatial side, to the constant
background $b/2$ appearing in the inverse transform of eq.~13's delta term
(main text §2.3), which is trivially smooth (a constant has infinite radius
of convergence everywhere) and carries no discontinuity. The pole $ib/k$
plays a complementary, quantitative role. Using the main text's globally
correct closed form $u(x,t)=b/2-(b/\pi)\arctan(x/(at))$ (the analytic
continuation of eq.~15--16 to all $x$) and the large-argument expansion
$\arctan(z)=\pi/2-z^{-1}+O(z^{-3})$,
\begin{equation}
u(x,t) - \frac{b}{2} \;\sim\; \frac{abt}{\pi x} \qquad (x\to+\infty,\ t\ \text{fixed}),
\label{frac:farfield}
\end{equation}
confirmed numerically: the ratio of the two sides is $0.996$ at $x=10$,
$0.999996$ at $x=100$, and $1-4\times10^{-7}$ at $x=1000$ (for
$a=1.3,b=1,t=0.8$). The algebraic, $1/x$ decay in~(\ref{frac:farfield}) is
the spatial signature of a simple pole at $k=0$: low-frequency
singularities in the spectrum correspond to algebraic, not exponential,
far-field tails. Neither survivor at $k=0$, the delta or the pole,
threatens the analyticity bound~(\ref{frac:stripwidth}), which is
controlled entirely by the decay of the multiplier at $|k|\to\infty$
established in~(\ref{frac:supbound}) and is completely insensitive to
whatever happens at $k=0$. The semigroup construction shows this smoothing
is not a special feature of one particular closed form: it is the generic
behaviour of $(-\Delta)^{1/2}$-semigroups, quantified
by~(\ref{frac:supbound})--(\ref{frac:stripwidth}) for any initial data,
independent of whatever singularity the spectrum retains at $k=0$.

\section{Centerpiece: the Caffarelli--Silvestre extension}
\label{sec:extension}

\subsection{The half-Laplacian is the derivative of the Hilbert transform}
\label{sec:hilbertsymbol}

Define the Hilbert transform $\HT v(x) := \dfrac{1}{\pi}\,\mathrm{P.V.}\displaystyle\int_{-\infty}^\infty \dfrac{v(y)}{x-y}\,dy$,
whose Fourier symbol, under the convention $\hat v(k)=\int v(x)e^{-ikx}dx$
used throughout this annex and the main text, is $-i\,\sgn(k)$. Composing
with differentiation, whose symbol is $ik$,
\begin{equation}
\mathcal{F}\!\left\{\frac{d}{dx}\big[\HT v\big]\right\}(k)
= (ik)\cdot(-i\,\sgn(k))\,\hat v(k)
= k\,\sgn(k)\,\hat v(k) = |k|\,\hat v(k),
\end{equation}
which is the symbol of $(-\Delta)^{1/2}$. So, as operators on
Schwartz functions,
\begin{equation}
(-\Delta)^{1/2}v = \frac{d}{dx}\big[\HT v\big].
\label{frac:hilbert-symbol}
\end{equation}
Confirmed numerically via FFT in \texttt{verify\_hilbert\_symbol.jl}. This
identity is specific to one spatial dimension: it is what makes the
half-Laplacian, uniquely among fractional powers, expressible through a
classical, purely local-looking construction (a derivative of a boundary
value of a harmonic function), which is the content of the next subsection.

\subsection{The extension problem at $s=\tfrac12$}

\citet{daoud2022fractional} (their fifth definition, following
\citealp{caffarelli2007extension}) realize $(-\Delta)^s$ as a
Dirichlet-to-Neumann map. Let $U(x,y)$ solve
\begin{equation}
\mathrm{div}\big(y^\alpha\nabla U\big) = 0 \quad\text{in } \mathbb{R}\times(0,\infty),
\qquad U(x,0) = v(x),
\qquad \alpha = 1-2s \in (-1,1).
\label{frac:extprob}
\end{equation}
Then
\begin{equation}
(-\Delta)^s v(x) = -d_s\lim_{y\to0^+} y^\alpha\,\partial_y U(x,y),
\qquad
d_s = \frac{2^{2s-1}\Gamma(s)}{\Gamma(1-s)}.
\label{frac:extformula}
\end{equation}
At $s=\tfrac12$: $\alpha = 1-2s = 0$ and
$d_{1/2}=2^0\,\Gamma(\tfrac12)/\Gamma(\tfrac12)=1$. The weight $y^\alpha$
disappears entirely, and~(\ref{frac:extprob}) collapses to the classical,
unweighted Laplace equation,
\begin{equation}
\boxed{\Delta U = 0 \quad\text{in } \mathbb{R}\times(0,\infty), \qquad
U(x,0)=v(x)}, \qquad
(-\Delta)^{1/2}v(x) = -\lim_{y\to0^+}\partial_y U(x,y).
\label{frac:extension}
\end{equation}
This is not a special case chosen for convenience: it is the \emph{only}
value of $s\in(0,1)$ for which the extension problem is the ordinary
harmonic-extension problem in the upper half-plane, with no singular or
vanishing weight at $y=0$.

For $N=1$, $U$ is the classical Poisson extension of $v$ into the upper
half-plane, explicitly
\begin{equation}
U(x,y) = \frac{y}{\pi}\int_{-\infty}^{\infty}
  \frac{v(x')}{(x-x')^2+y^2}\,dx',
\label{frac:poisson}
\end{equation}
recovering $v$ as $y\to0^+$ and solving $\Delta U=0$ throughout $y>0$; and
its harmonic conjugate $V$ (so that $U+iV$ is holomorphic in
$y>0$) has boundary trace equal to the Hilbert transform of $v$,
$V(x,0)=\HT v(x)$, a classical fact of harmonic and complex analysis. The
Cauchy--Riemann relations give $-\partial_yU = \partial_xV$, so at $y=0$,
\begin{equation}
-\lim_{y\to0^+}\partial_yU(x,y) = \frac{d}{dx}\big[V(x,0)\big]
= \frac{d}{dx}\big[\HT v\big](x),
\end{equation}
which is $(-\Delta)^{1/2}v(x)$ by~(\ref{frac:hilbert-symbol}),
confirming~(\ref{frac:extension}) is internally consistent with the
independent Fourier-symbol derivation of Section~\ref{sec:hilbertsymbol}.

\subsection{Physical reading}

Equation~(\ref{frac:extension}) says that solving for $(-\Delta)^{1/2}$ at
the surface of $\mathbb{R}$ is the same problem as solving Laplace's
equation in the half-space $y>0$ above it and reading off the normal
derivative at the boundary. That is, structurally, the same type of problem
Savage's Stokes kernel already comes from: a linear response at the surface
of a half-space, generated by potential theory in the half-space beneath it.
Two-dimensional half-space Green's functions, whether for an elastic solid
(the Boussinesq--Flamant problem) or for Stokes flow (via the Goursat
representation of the biharmonic stream function as a pair of harmonic
functions, the standard reduction used throughout half-space viscous flow
theory), are built from this kind of harmonic/conjugate-harmonic
machinery. This is a structural argument about why half-space linear
response problems in two dimensions generically produce Hilbert-transform-type,
half-Laplacian operators; it is not a claim to have reproduced
\citeauthor{savage1978asthenosphere}'s (\citeyear{savage1978asthenosphere})
original derivation step by step. What Section~\ref{sec:kernel} establishes
independently, by direct computation on the kernel actually used in the main
text, is the concrete fact: whatever the exact route through Stokes flow
theory, the resulting kernel is, exactly and with the correct constant and
sign, $-2\eta$ times $(-\Delta)^{1/2}$. Equation~(\ref{frac:extension})
explains why an operator of that specific type, half-order and no other, is
the one a linear two-dimensional half-space coupling should be expected to
produce.

\subsection{Making the analogy precise: the true field is biharmonic, not harmonic}

The structural claim above can be sharpened into an explicit computation.
Savage's kernel is the response of an actual two-dimensional Stokes flow,
not of the abstract harmonic extension~(\ref{frac:extension}), and the two
problems solve different equations. Writing $\zeta:=z-H$ for the
depth below the crust--mantle interface (main text's $z$ increases
downward, with the interface at $z=H$), incompressible Stokes flow in
$\zeta>0$ is governed, through the stream function $\psi$
($v_x=\partial_\zeta\psi$, $v_z=-\partial_x\psi$, so that incompressibility
holds automatically), by the biharmonic equation $\nabla^4\psi=0$, not by
Laplace's equation.

Solving it directly confirms, and sharpens, the structural analogy.
Fourier-transforming in $x$, the biharmonic equation reduces to the
ordinary differential equation $(\partial_\zeta^2-k^2)^2\hat\psi=0$, whose
solution bounded as $\zeta\to\infty$ is
$\hat\psi(k,\zeta)=(A+B\zeta)e^{-|k|\zeta}$. Imposing no vertical flow
through the interface, $\hat\psi(k,0)=0$, fixes $A=0$; imposing the
prescribed surface velocity, $\partial_\zeta\hat\psi(k,0)=\hat V(k,t)$
where $V:=\partial u/\partial t$, fixes $B=\hat V(k,t)$:
\begin{equation}
\hat\psi(k,\zeta,t) = \hat V(k,t)\,\zeta\,e^{-|k|\zeta}.
\label{frac:stokesdepth}
\end{equation}
This gives the velocity and stress fields at any depth $\zeta$ in the
asthenosphere, not only at $\zeta=0$. Taking the shear stress
$\tau_{xz}=\eta(\partial_\zeta^2\psi-\partial_x^2\psi)$ at the interface
reproduces Savage's kernel exactly:
$\hat\tau_{xz}(k,0,t) = -2\eta|k|\hat V(k,t)$, matching~(\ref{frac:kernel-identity}).

The mismatch with~(\ref{frac:extension}) is now explicit rather than merely
asserted. A pure harmonic extension of the same boundary data would decay
as $e^{-|k|\zeta}$ alone; the physical Stokes solution~(\ref{frac:stokesdepth})
carries an extra factor of $\zeta$, the signature of the repeated root at
rate $|k|$ in the biharmonic ODE, and the structure Goursat's
representation predicts (a harmonic function multiplied by a coordinate,
combined with a second harmonic function). So the physical half-space
field is not literally the harmonic extension of
Section~\ref{sec:extension}; it is a two-harmonic-function, biharmonic
object that happens to produce the same boundary operator,
$-2\eta(-\Delta)^{1/2}$, once the interface stress is extracted. The
half-Laplacian is the correct description of the boundary response either
way; the extension variable $y$ of~(\ref{frac:extprob}) should not be read
as the physical depth $\zeta$, and depth-resolved fields in the true
asthenosphere require solving~(\ref{frac:stokesdepth}) directly, not
substituting the surface solution into the Poisson kernel~(\ref{frac:poisson}).

\section{Proofs}
\label{sec:proofs}

The short proofs load-bearing for this model's own construction are given
inline, at the point of use, in Sections~\ref{sec:kernel}--\ref{sec:extension}
(the value of $C(1,\tfrac12)$, the P.V./Hadamard sign identity, the Sobolev
divergence of the step function, the Hilbert-transform symbol identity, and
the $\alpha=0$ reduction of the extension problem). This section collects
the two longer, more general results used along the way, for which we cite
rather than reproduce the full argument, following the same practice as
\citet{daoud2022fractional} themselves.

\paragraph*{Equivalence of the singular-integral and Fourier definitions,
general $N$.} That~(\ref{frac:singint}) and the Fourier-multiplier
definition $(-\Delta)^sv = \mathcal{F}^{-1}\{\|k\|^{2s}\hat v(k)\}$ agree for
every $N\ge1$ and $s\in(0,1)$, not only the $N=1$, $s=\tfrac12$ case proved
directly in Section~\ref{sec:kernel}, is Definitions~2.1 and~2.2 in
\citet{daoud2022fractional}, with the full proof referred there to
\citet[Section~3]{dinezza2012hitchhikers}. The general argument decomposes
the singular integral via a change to polar coordinates and integrates the
angular part in closed form against the Fourier representation; it is not
reproduced here because it is not needed beyond the $N=1$ case established
directly above.

\paragraph*{Full derivation of the extension formula.} The extension
problem~(\ref{frac:extprob})--(\ref{frac:extformula}) for general
$s\in(0,1)$ is proved in \citet{caffarelli2007extension}, via an explicit
Poisson-type kernel for the weighted operator $\mathrm{div}(y^\alpha\nabla\cdot)$
and a direct computation of its Fourier transform in $x$ for fixed $y$,
showing the resulting ODE in $y$ has the modified Bessel function solution
whose $y\to0^+$ weighted normal derivative reproduces $\|k\|^{2s}\hat v(k)$
up to the constant $d_s$. We used only the $\alpha=0$ specialization, proved
directly and elementarily in Section~\ref{sec:extension} via the classical
Poisson-kernel/harmonic-conjugate construction, which does not require the
general Bessel-function machinery.

\section{Synthesis for structural geology}
\label{sec:synthesis}

The mathematics above answers three questions a structural geologist would
otherwise have to take on faith from the main text.

Why does deformation appear everywhere at once, with no propagating front?
Because the operator is genuinely non-local: Section~\ref{sec:kernel} shows
the Stokes kernel is exactly $-2\eta$ times $(-\Delta)^{1/2}$, an operator
whose value at any point $x$ depends on the field at every other point
simultaneously, not on a neighborhood that grows with time as in ordinary
diffusion. This is the same conclusion the main text reaches in §4.3 by
comparing Lorentzian and erfc tails; here it follows directly from the
kernel's identity as a fractional Laplacian rather than from the shape of
the fitted curves.

Why does an instantaneous dike opening not break the model? Because the
half-Laplacian semigroup is smoothing, and quantifiably so:
Section~\ref{sec:semigroup} shows that no matter how singular the initial
displacement field is (a bare Heaviside step fails to belong to the
model's own natural energy space, Section~\ref{sec:heaviside}), the field
becomes real analytic in space for any $t>0$ after the intrusion, with a
radius of convergence that grows linearly, at least $at$, at every point.
The arctan step response is not a lucky closed form; it is the generic
behaviour of any field evolving under this operator. This regularization is
intrinsic to the operator, not to the data: the radius $at$ comes entirely
from how fast the semigroup multiplier $e^{-a|k|t}$ decays at high
frequency, a property that holds for any initial condition in the space of
tempered distributions, however singular; the dike-opening step is one
instance of that generic behaviour, not something arranged for it in
particular. The same section gives a second, complementary picture of why:
the correct solution can be
built by averaging the ordinary (Fickian) diffusion response to the same
initial condition over every possible diffusion timescale, weighted by a
fixed subordinator density. That ordinary diffusion response is the Foulger
profile the main text rejects as the physical model (§4.3); it
survives here only as a mathematical ingredient, blended across timescales
rather than taken at face value at any one of them.

Why is the operator order exactly one half, and not some other exponent
tuned to fit the data? Because, as shown in Section~\ref{sec:extension},
$s=\tfrac12$ is the unique order at which the extension problem defining
$(-\Delta)^s$ reduces to the ordinary harmonic problem in the half-space
beneath the surface. This is not only a structural analogy: the same
section confirms it by direct computation, solving the actual
two-dimensional Stokes flow beneath the plate (governed by the biharmonic
equation, not the harmonic one) and showing its interface stress
reproduces $-2\eta(-\Delta)^{1/2}$ exactly, with the extension problem's
variable $y$ standing in for, but not equal to, the true physical depth.
A Newtonian asthenosphere coupled to a thin elastic plate could not have
produced any other operator order: half-order non-locality is the
structural signature of this specific rheology and this specific
dimensionality, not a fitted parameter.

That last point is also where the mathematics opens a direction for
future work rather than closing the subject. The derivation is tied
throughout to Newtonian viscosity; a power-law (non-Newtonian) asthenosphere
of the kind independently documented in the mantle beneath magmatically
active rifts would replace the linear Stokes kernel with a different
constitutive relation, and there is no guarantee its Fourier reduction
retains the clean $|k|$ dependence exploited here. Whether such a rheology
still reduces to a fractional Laplacian of some order $s\neq\tfrac12$, or to
an operator outside this family altogether, is an open question this annex
does not attempt to answer, but the extension-problem machinery of
Section~\ref{sec:extension} gives a concrete route to it: any proposed
kernel could be tested against the same $\alpha=1-2s$ reduction to see
whether a consistent order $s$ exists at all.

\bibliographystyle{plainnat}
\bibliography{../paper/refs}

\end{document}
